\clearpage
\fancyhead[L]{\small CONNAISSANCE DES TEMS.}
\hypertarget{native-643-lower}{}
\begin{center}\small Göttingische gelehrte Anzeigen. Stück 80. Seite 793..796. 1824 Mai 17.\\[3pt]\rule{.72\linewidth}{.3pt}\end{center}

\emph{Connaissance des tems ou des mouvemens célestes à l'usage des astronomes et des navigateurs pour l'an \textup{1826}. Publiée par le bureau des longitudes.} Paris 1823. Bei Bachelier. Die stehenden Artikel 216 Seiten, die Zusätze 110 Seiten. gr. 8.

Zie Zusätze zu diesem Jahrgange bestehen in folgenden Artikeln. Uebersicht der astronomischen Beobachtungen, die in den Jahren 1820..1822 auf der Marseiller Sternwarte angestellt sind von \textsc{Gambart} d.\,J. Die Marseiller Sternwarte ist bisher mit wenigen und mittelmässigen Instrumenten versehen gewesen, und Hr. \textsc{Gambart} hat sich daher mit Recht auf solche Beobachtungen beschränkt, die diesem Umstande angemessen sind. Der fast immer heitere Himmel von Marseille ist besonders einladend zum Aufsuchen neuer Cometen, und das Glück, welches dabei früher den dadurch berühmt gewordenen \textsc{Pons} so sehr begünstigt hat, scheint auch \textsc{Gambart} treu zu bleiben, der von den drei Cometen des Jahrs 1822 einen ganz zuerst, und einen wenigstens ohne von der frühern Entdeckung in Marlia zu wissen, selbst entdeckt hat. \textsc{Gambart}'s schätzbare Beobachtungen dieser beiden Cometen sind hier mit aller Ausführlichkeit, die man wünschen konnte, abgedruckt. Auch an dem zweiten Cometen des Jahres 1822 hat \textsc{Gambart} zwei Beobachtungen gemacht, die um so schätzbarer sind, da dieser Comet übrigens nur noch von \textsc{Caturegli} in Bologna beobachtet ist. Die übrigen Beobachtungen bestehen in Verfinsterungen von Jupiterstrabanten, Sternbedeckungen und der Mondsfinsterniss vom 2. August 1822. --- Ueber die Vertheilung der Wärme in einem homogenen Ringe, dessen Dicke überall dieselbe ist, und der sich in einem Raume befindet, in dessen verschiedenen Punkten ungleiche Temperatur Statt hat [von \textsc{Poisson}]. Obgleich dieser schöne Aufsatz, wegen der feinen mathematischen Behandlung, hauptsächlich den Geometer interessirt, so ist er doch auch nicht ohne Interesse für die praktische Astronomie, insofern Kreis-Instrumente sich beim Beobachten nicht selten in solchen Umständen befinden, wie hier zum Grunde gelegt \SourceLeaf{644} sind, und es wünschenswerth ist, von der daraus entstehenden Defiguration der Instrumente bestimmte Vorstellungen zu erhalten. --- Ueber die Geschwindigkeit des Schalles [von \textsc{Poisson}]; eine eben so gründliche als lichtvolle Darstellung der durch \textsc{Laplace} zuerst vollendeten Theorie dieses Gegenstandes: die Benutzung des Verhältnisses der specifischen Wärme der Luft bei gleichbleibendem Druck zu der specifischen Wärme bei gleichbleibendem Volumen, wie solches aus dem interessanten Versuche von \textsc{Clement} und \textsc{Delorme} hervorgeht, bringt die Theorie der Geschwindigkeit des Schalles mit den Resultaten der Erfahrung in beinahe vollkommene Uebereinstimmung. --- Beobachtungen des ersten und dritten Cometen vom Jahr 1822 auf der Pariser Sternwarte, der letztere war von Hrn. \textsc{Bouvard} selbst, obwohl einige Tage später als von \textsc{Pons} und \textsc{Gambart}, entdeckt. --- Resultate aus Beobachtungen mit unveränderlichen Pendeln berechnet von Hrn. \textsc{Mathieu}. Der Schiffslieutenant \textsc{Duperrey}, welcher in den Jahren 1822 und 1823 mit der Corvette Lacoquille eine Reise um die Welt machte, führte zwei unveränderliche Pendel mit sich, deren Schwingungen vor der Abreise in Paris, und von Hrn. \textsc{Duperrey} bei einer Landung auf den Malouinen beobachtet wurden. Da die südliche Breite der letztern nur $2^\circ\,41'$ grösser ist, als die nördliche Breite von Paris, so lässt sich aus diesen Versuchen, die übrigens mit vieler Sorgfalt discutirt sind, eigentlich nur das folgern, dass sie mit einer Abplattung von $\frac{1}{308}$ nicht im Widerspruch sind. Wichtiger sind die hier gleichfalls berechneten Versuche mit einem unveränderlichen Pendel, welche \textsc{Brisbane} und \textsc{Rümker} zuerst in London und nachher in Paramatta auf Neuholland angestellt haben, und aus denen hier die Abplattung $\frac{1}{206}$ berechnet wird. --- Ueber die Wirkung des Mondes auf die Atmosphäre, von Hrn. \textsc{de Laplace}. Seit dem Jahr 1815 wird auf der Pariser Sternwarte die Barometerhöhe täglich viermal, nemlich Vormittags um 9 Uhr, Mittags, Nachmittags 3 Uhr und 9 Uhr beobachtet. Die tägliche Variation des Barometerstandes tritt aus den zufälligen Schwankungen doch so bestimmt hervor, dass der Unterschied zwischen dem Stande Vormittags 9 Uhr und Nachmittags 3 Uhr, nach seinem Mittelwerthe für jeden einzelnen Monat während sechsjähriger Beobachtungen auch nicht ein einzigesmal sein Zeichen geändert hat (im Mittel ist dieser Unterschied 0.801 Millimeter, und zwar die vormittägige Höhe die grössere). Dies veranlasste den Versuch, aus einer zweckmässigen Auswahl dieser Beobachtungen auch den Einfluss des Mondes auf den Barometerstand auszumitteln. Es wurden, zu diesem Zwecke, die Beobachtungen in der Nähe der Syzygien mit denen in der Nähe der Quadraturen verglichen; das Resultat dieser Vergleichung war eine kleine periodische Ungleichheit, deren Periode ein halber Mondstag ist, und deren Maximum nur 0.0272 Millimeter beträgt, und am Tage der Syzygien um $3^{\rm u}18^{\rm m}36^{\rm s}$ Statt findet. Da in den numerischen Theil dieser Untersuchung einige Unrichtigkeiten eingeschlichen zu sein scheinen, so wurde Ref. hiedurch und durch das Interesse des Gegenstandes veranlasst, diese Rechnungen theils schärfer, theils nach andern Methoden zu wiederholen. Das Resultat ist aber fast dasselbe, wie das von \textsc{Laplace} erhaltene; Ref. findet nemlich das Maximum 0.0269 Millimeter, dessen Zeit $3^{\rm u}11^{\rm m}49^{\rm s}$, und den mittlern bei der ersten Bestimmung zu befürchtenden Fehler 0.0264 Millimeter. --- Den Beschluss dieses Jahrganges macht ein kleiner Aufsatz von Hrn. \textsc{Puissant} über die Berechnung der mit einem Repetitions-Theodolithen gemessenen Azimuthe. Das hier vorgeschlagene Verfahren ist nicht neu, sondern im wesentlichen ganz einerlei mit dem, welches Hr. \textsc{Soldner} in den Denkschriften der Münchner Akademie für 1813 abgehandelt hat, und worüber auch unsre Blätter Jahrg. 1815 S. 449 ff. [S. 587 d. B.] nachgesehen werden können.
